\documentclass[landscape,twocolumn,leqno,a4paper,10pt]{report} %leqno affiche les numéros d'équations à gauche

\usepackage[a4paper,top=0.4cm,bottom=1cm,left=0.4cm,right=0.4cm,marginparwidth=0.5cm,footskip=0.6cm]{geometry}



%----------------- Codage de la langue et de la police--------------------------------------------------------------------------
\usepackage[french]{babel} 		% Permet de spécifier au compilateur que l'on désire taper son document en français.
\usepackage[utf8x]{inputenc} 	% Permet de taper directement à l'écran, dans le fichier .tex, les caractères accentués.
\usepackage[T1]{fontenc} 		% Permet juste d'utiliser une nouvelle norme LATEX concernant le codage des caractères.
%\usepackage{lmodern}		% Pack de police --> Pas besoin puisque j'utilise Fourier

%------------------------ Packages utiles-------------------------------------------------------------------------------------------

\usepackage{graphicx}	% Pour insérer des images
\usepackage{caption}		% Légende

\newcommand\hmmax{0}		% Alors c'est une astuce assez mystérieuse pour dépasser la limite de polices latex.
							% En effet, le package bm est lourd et avec tout ce que l'on met, on obtient un message d'erreur : Too many math alphabets used in version normal.			
\newcommand\bmmax{0}		% Cette astuce règle le problème, il faut tout de même mettre ces compteurs avant le package bm

\usepackage{amsmath,amssymb}	% Permettent d'obtenir des caractères mathématiques (de toute beauté !).
\usepackage{amsthm}		% Pour pouvoir faire les newthm
\usepackage{mathtools}	% Corrige des bugs dans amsmath et le complète
\usepackage{numprint}	% Pour les espacements dans les nombres : \numprint{10000} = 10 000
\usepackage{fourier}		% REMARQUE : [upright] impose aux majuscules d'être droite dans le mode mathématique 
\usepackage{pifont} 		% Importe une large gamme de symboles, on utilise \ding{numéro}
\usepackage{xargs} 		% Permet de définir des macros à plusieurs paramètres optionnels.
\usepackage{fancyhdr} 	% Package pour les entêtes et bas de page
\pagestyle{fancy} 			% Accompagne \usepackage{fancyhdr}
\usepackage{bm} 		% Permet de mettre en gras du contenu mathématiques via \bm{..}
\usepackage{mathrsfs} 	% Pour les lettres cursives
\usepackage{enumitem} 	% Pour pouvoir personnaliser les listes

\frenchbsetup{StandardLists=true} % Afin d'éviter tout conflit entre les extensions enumitem et french babel. 

\usepackage{sectsty}	% Pour colorer les titres
\usepackage{titlesec} 	% Pour personnaliser les titres
\usepackage{setspace} 	% Pour changer la taille des interlignes
       	               
\usepackage{array,multirow,makecell}
\usepackage{hlist}		% Package Christian Tellechea sur les listes numérotées gauche vers droite
\setdefaulthlist{left margin = 2mm, label sep= 0.2mm}	% réglages package hlist
				

%------------ Les packages de dessin sous pstricks : il faut changer la compilation de latex pour compiler ces dessins------------------
% Options\Configurer Texmaker\Compil rapide\ Latex+dvips+ps2pdf+Voir PDF

\usepackage[pdf]{pstricks} 	% Pour faire des dessins
\usepackage{pstricks-add,pst-plot,pst-tree,pst-node,pst-math,pst-eucl,pst-func,pst-3dplot}% ,pst-text,pst-eps,pst-fill,pst-grad
\usepackage{multido}		% Pour faire des multido
\usepackage{lastpage} 	% Pour le compteur de la dernière page, on appelle alors celle-ci avec lastage



\usepackage[colorlinks=true, 
			%allcolors=blue
			]{hyperref}		% ATTENTION : doit être importé après la majorité des packages sauf tabularx et breakrules (par exemple)
\hypersetup{			% Paramétrage des hyperliens
	colorlinks = true,   		% Colorise les liens
	breaklinks = true,		% Permet les retours à la ligne pour les liens trop longs
	urlcolor = blue,		% Couleur des hyperliens
	linkcolor = black,		% Couleur des liens internes aux documents (index, figures, tableaux, équations, ...)
	citecolor = blue		% Couleur des liens vers les références bibliographiques
	}




%%%%%%%%%%%%%%%%%%%%%%%%%%%%%%%%%%%%%%%%%%%%%%%%%%
%------------------------- Définition des commandes -----------------------------------------------------------------------------
%%%%%%%%%%%%%%%%%%%%%%%%%%%%%%%%%%%%%%%%%%%%%%%%%%


%--------------------- Raccourcis fontes--------------------------------------------------------------
% REMARQUE : le package fourier met la même police pour mathcal et mathscr, on remet donc mathcal normalement et on utilise la définition de fourier (un peu plus grasse) pour mathscr
\DeclareMathAlphabet{\mathcal}{OMS}{cmsy}{m}{n}	
\DeclareMathAlphabet{\mathscr}{FMS}{futm}{m}{n}

%---------------------------- Raccourcis des espaces--------------------------------------------------------------------------------------------
\newcommand\n{\noindent}
\newcommand\q{\quad}
\newcommand\qq{\quad\quad}

%--------------------- Raccourcis en mode mathématiques-------------------------------------------------------------------------------------
\newcommand{\et}{\text{ et }}
\newcommand{\ou}{\text{ ou }}
\newcommand{\SSI}{\Longleftrightarrow}
\let\si\relax
\newcommand{\si}{\text{ si }}
\newcommand{\sinon}{\text{ sinon }}
\renewcommand{\epsilon}{\varepsilon}		% epsilon beaucoup plus joli
\renewcommand{\geq}{\geqslant} 
\renewcommand{\leq}{\leqslant}
\newcommand{\e}{\mathrm{e}}
\renewcommand{\textbf}[1]{\begingroup\bfseries\mathversion{bold}#1\endgroup}		% Très important, permet de passer en gras même avec le mode mathématiques

\newcommand{\Mod}[1]{\allowbreak\if@display\mkern5mu\else\mkern5mu\fi\left[#1\right]}	% Pour modulo entre crochets [.] avec un espacement convenable 


%---------------------------------- Raccourcis calcul diff----------------------------------------------------------------------------------------
\renewcommand{\d}{\;\mathrm{d}}	% Pour le d droit de l'intégrale car ce n'est pas une variable...


%-------------------------------- Raccourcis des ensembles----------------------------------------------------------------------------
\newcommand\R{\mathbb{R}}
\newcommand\C{\mathbb{C}}
\newcommand\K{\mathbb{K}}
\newcommand\Q{\mathbb{Q}}
\newcommand\D{\mathbb{D}}
\newcommand\Z{\mathbb{Z}}
\newcommand\N{\mathbb{N}}
\newcommand\U{\mathbb{U}}


%------------------------------- Dimensions intelignes, items...-----------------------------------------------------------------------------------------------
\onehalfspacing	% Pour que l'interligne soit de 1.5. Pour que l'interligne soit de la valeur que l'on souhaite, de base : \setstretch{1}

\setlist[1]{font=\bfseries,itemsep=0.2mm,leftmargin=6mm}
\setlist[2]{font=\bfseries,itemsep=0.2mm,leftmargin=4mm}
\setenumerate[1]{font=\bfseries,itemsep=0.2mm,leftmargin=6mm}
\setenumerate[2]{font=\bfseries, label=\alph*., leftmargin=4mm}
%%\setenumerate[profondeur]{font=\bfseries, label=\ding{221} \color{blue}} permet de changer en a) b) les sous-items d'une liste enumerate
%% itemsep = correspond à l'interligne entre chaque item
% de base : labelsep=2mm et permet de changer la distance entre l’étiquette et le texte des éléments
% leftmargin =4mm spécifie la marge à gauche de la liste
\AtBeginDocument{\renewcommand{\labelitemi}{\textbullet}}	% Pour que les puces par défaut soit des points.
%\AtBeginDocument{\renewcommand{\labelitemi}{\ding{86}}}	% Pour que les puces par défaut soit des étoiles.
\AtBeginDocument{\renewcommand{\labelitemii}{$\ast$}}		% Pour que les puces par défaut pour une sous-liste soit des astérisques.
\newcommand{\itemmm}{\item[$\ast$]}				% \itemmm devient l'item \ast
\newcommand{\itemmmm}{\item[\footnotesize$\circledcirc$\normalsize]}
\newcommand{\itemmmmm}{\item[\footnotesize$\astrosun$\normalsize]}
\newcommand{\itemmmmmm}{\item[\footnotesize$\multimap$\normalsize]}
\newcommand{\itemmmmmmmmm}{\item[$\divideontimes$]}
\newcommand{\etoile}{\item[$\star$]}
\newcommand{\etoilee}{\item[$\bigstar$]}
\newcommand{\losange}{\item[$\diamondsuit$]}
\newcommand{\carre}{\item[\LARGE$\centerdot$]}
\newcommand{\carreombre}{\item[\ding{112}]}

\setlength{\columnsep}{30pt} 		% Permet d'écarter les deux colonnes du twocolumn de 30pt, pratique pour écarter le texte de la ligne centrale
\setlength{\columnseprule}{0.5pt}	% Affiche une ligne verticale entre chaque item et entre chaque page du twocolumn
						% \columnseprule=0cm 	% Permet d'enlever le trait vertical entre chaque colonne
%\renewcommand\tabenumsep{0.5cm\hskiplem}	% Ecarte de 0.5cm les items de tabenum, sinon ils sont trop collés...


\newlength{\tagmarginsep} 		% On crée une longueur
\setlength{\tagmarginsep}{0cm}		% Permet de régler la position d'un tag pour une équation mettre avant l'équation : \setlength{\tagmarginsep}{x cm} pour que le tag se rapproche de l'équation
\everydisplay{\displayindent=\tagmarginsep \displaywidth=\dimexpr\linewidth-2\tagmarginsep}	% Va avec ce qu'il y a juste avant


%-------------------- Définition des couleurs-----------------------------------------------------------------------------------------

\definecolor{violetfonce}{RGB}{126,62,142}
\definecolor{magentafonce}{RGB}{128, 0, 128}
\newrgbcolor{magentafonce}{.5 0 .5}			% Une autre façon de définir une couleur, permet surtout d'utiliser \magenta
\definecolor{violet1}{cmyk}{0.5,1,0,0}
\definecolor{violet2}{cmyk}{0.24998, 1, 0, 0}
\definecolor{violet3}{cmyk}{0, 1, 0.25, 0}



\newcommand\tred{\textcolor{red}}
\newcommand\tredf{\textcolor{rougefonce}}
\newcommand\tdred{\textcolor{deepred}}
\newcommand\tblue{\textcolor{blue}}
\newcommand\tbluef{\textcolor{bleufonce}}
\newcommand\tgreen{\textcolor{vertfonce}}
\newcommand\torange{\textcolor{orange}}
\newcommand\tmagenta{\textcolor{magentafonce}}


%------------------ Changement paramètres des parties, sous parties---------------------------------------------------------


\sectionfont{\color{violet1}}			% Utilisation du package sectsty
\subsectionfont{\color{violet2}}
\subsubsectionfont{\color{violet3}}


%--------------------------------- Raccourcis probabilités--------------------------------------------------------------------------------
\newcommand{\Prob}{\mathbb{P}}
\newcommand{\E}{\mathbb{E}}
\newcommand{\normale}[2]{\mathcal{N}\left(#1;#2 \right)}	% Pour la loi Normale : 	$\normale{\mu}{\sigma^2}$
\newcommand{\Unif}[1]{\mathcal{U}\left(#1\right)}			% Pour la loi uniforme :	$\Unif{[a;b]}$
\newcommand{\Ber}[1]{\mathcal{B}\left(#1\right)}			% Pour la loi de Bernoulli : $\Ber{p}$
\newcommand{\Bino}[2]{\mathcal{B}\left(#1;#2\right)}		% Pour la loi Binomiale :	$\Bino{n}{p}$
\newcommand*{\binomial}[2]{% 
  \ensuremath{ 
    \begin{pmatrix} 
      #2\\ 
      #1
    \end{pmatrix}}}				% Pour le coefficient binomial : $\binom{k}{n}$
\renewcommand{\binom}{\binomial} 	% La commande \binom existait déjà mais n'était pas pratique



%----------------------------------- Opérateurs---------------------------------------------------------------------------------------------------

\newcommand{\barre}[1]{\overline{#1}}  									% Barre au dessus du texte, pratique pour le conjugué des complexes
\newcommandx{\Sum}[3][3=\white]{{\displaystyle\sum\limits_{\substack{#1 #3}}^{#2}}}		% $\Sum{k=0}{n}[<option>]$ où [<option>] est placée à la suite de k=0 si l'on tape de façon classique, et en-dessous si l'on met [\\ k\neq k_0]
															% On peut juste taper : $\Sum{k=0}{n}$ et cela fonctionne
															% \displaystyle impose de gérer les tailles comme en mode mathématique : ici on aura donc toujours de grande somme que ce soit en math comme dans le texte
\newcommandx{\Prod}[3][3=\white]{{\displaystyle\prod\limits_{\substack{#1 #3}}^{#2}}}		% $\Prod{k=0}{n}[<option>]$	idem pour l'option
\newcommand{\Int}[4]{{\displaystyle\int_{#1}^{#2} #3  \d #4}}   					% $\Int{0}{\pi}{\cos(t)}{t}$
\newcommandx{\Lim}[3][3=\white]{\lim\limits_{\substack{#1 \to #2 #3}}}		% $\Lim{x}{0}[<option>]$	idem pour l'option


%------------------------------------ Délimiteurs--------------------------------------------------------------------------------------------------
%\newcommand{\abs}[1]{\left\lvert #1 \right\rvert}			% Valeur absolue | |
%\newcommand{\norm}[1]{\left\lVert #1 \right\rVert}			% Doubles barres verticales : norme d'un vecteur || ||

%--------------------------------------- Vecteurs----------------------------------------------------------------------------------------------------
 
\newcommand{\Oij}{\left(O;\vec{\imath},\vec{\jmath}\right)}
\newcommand{\Oijk}{\left(O;\vec{\imath},\vec{\jmath},\vec{k}\right)}
\newcommand{\Ouv}{\left(O;\vec{u},\vec{v}\right)}
\newcommand{\Ouvw}{\left(O;\vec{u},\vec{v},\vec{w}\right)}

%------------------------------------ Systèmes d'équations-------------------------------------------------------------------------------------------
% Remarque \syst et \syste ne font pas d'alignement mais ils sont bien, si l'on veut de l'alignement, on utilisera \systeme du package "systeme"
\newcommand*{\syst}[2]{%			% Systèmes 2 éq. à ? inconnues : 	\syst{ax+by+c=0}{a'x+b'y+c'=0}
  \left\{\begin{array}{ll} 
        #1\\ 
        #2 
      \end{array}  
    \right. }
 
\newcommand*{\syste}[3]{% 			% Systèmes 3 éq. à ? inconnues : 	\syste{ax+by+c=0}{a'x+b'y+c'=0}{a''x+b''y+c''=0}
  \left\{\begin{array}{ll} 
        #1\\ 
        #2\\
        #3 
      \end{array}  
    \right.}



%------------------------------------------ Autre----------------------------------------------------------------------------------------------------------
\newcommand{\boite}[2]{\noindent\parbox[]{#1 cm}{#2} }  									% Boite autre façon de faire minipage en gros


%---------------------------- En-têtes et bas de pages----------------------------------------------------------------------------
% A chaque fois : *head* pour le haut et *foot* pour le bas
% L = left, C = center, R = right 
% L = toutes les pages, LE = toutes les pages paires, LO = toutes les pages impaires, idem CE, CO, ...


\newcommand{\fcours}[1]{								% f pour fancy				   
	\pagestyle{fancy}									% Efface le style qui a pu être mis avant
	\renewcommand{\headrulewidth}{0pt} 
	\renewcommand{\footrulewidth}{0.5pt} 
	%\fancyhead[L]{}									
	%\fancyhead[C]{}
	%\fancyhead[R]{}									
	\fancyfoot[L]{\vspace{-1mm}\footnotesize \textit{#1}} 			% Classe : <Seconde> 
	\fancyfoot[C]{\vspace{-1mm}\footnotesize \thepage/\pageref{LastPage}}
	\fancyfoot[R]{\vspace{-1.8mm}\textit{\footnotesize C. Marguerite}}	% N'apparaîtra pas si on met les marges avec moins de bottom=1.3cm dans le package geometry
 }												% Fin de la commande
					%% Mettre au début du document : \fcours{Seconde} par exemple
					%% \usepackage[a4paper,top=0.4cm,bottom=1cm,left=0.4cm,right=0.4cm,marginparwidth=0.5cm,footskip=0.6cm]{geometry} 




%---------------------------- Réglages chapitres + table des matières------------------------------------------------------
\titleformat{\part}[display]
{\centering\normalfont\Huge\bfseries}
{\partname}
{20pt}
{\rule{6cm}{1pt}\raisebox{-2.1pt}[10pt][10pt]{\quad\decofourleft\decoone\decofourright\quad}\rule{6cm}{1pt} ~\\\vspace{20pt}}

\makeatletter
\renewcommand\chapter{\global\@topnum\z@
                    \@afterindentfalse
                    \secdef\@chapter\@schapter}
\def\@topnewpage[#1]{#1}
\makeatother										% Pour avoir le chapitre sur une demie page

\titleformat{\chapter}{\centering\normalfont\Large\bfseries\scshape\color{magentafonce}}{}{1em}{}
\titlespacing*{\chapter}{0pt}{-18pt plus 0pt minus 2pt}{0pt plus 2pt minus 2pt}		% Réglages espacements (à gauche, au-dessus, en-dessous) de la commande \chapter

\makeatletter
\newcommand{\twocolumntoc}{		
  \section*{\huge\centering\contentsname
    \@mkboth{%
      \MakeUppercase\contentsname}{\MakeUppercase\contentsname}}%
  \@starttoc{toc}%
}
\makeatother										% Permet de mettre la table des matières en twocolumn 
												% Mettre après \begin{doc..} \twocolumntoc \clearpage 		% Remplace \tableofcontents par \twocolumntoc en twocolumn

\makeatletter
\renewcommand{\thesection}{\@arabic\c@section}
\makeatother										% Enlève la numérotation des chapitres devant les sections

\setcounter{tocdepth}{3} 								% Permet d'afficher dans la table des matières suivant une certaine profondeur : -1=part, 0=chapter, 1=section, ...
\setcounter{secnumdepth}{3}							% Pour avoir la numérotation des subsubsections avec les classes book et report

\newcommand{\nocontentsline}[3]{}
\let\oldaddcontentsline\addcontentsline
\newcommand{\tocless}[2]{\let\addcontentsline\nocontentsline #1{#2}\let\addcontentsline\oldaddcontentsline}	% Pour retirer une section de la table des matières : \tocless\section{...}





\setstretch{1}
\begin{document}
\fcours{Concours général}

\chapter{Concours général des lycées - Session 2004}

\section{Une famille de cercles tangents}
 
\n Dans le plan, soit $A$ un point et $\Delta$ une demi-droite d'origine $A$.

\begin{enumerate}
	\item On considère trois cercles $\Gamma_{1}$, $\Gamma_{2}$, $\Gamma_{3}$ de même rayon $r$ non nul, de centres respectifs $O_{1}$, $O_{2}$, $O_{3}$ distincts et alignés dans cet ordre sur la demi-droite $\Delta$. Le cercle $\Gamma_{1}$ passe par $A$ et le cercle $\Gamma_{2}$ est tangent aux cercles $\Gamma_{1}$ et $\Gamma_{3}$.\\
	Les diamètres des cercles $\Gamma_{1}$, $\Gamma_{2}$, $\Gamma_{3}$ sur la demi-droite $\Delta$ sont notés respectivement $\left[A A_{1}\right]$, $\left[A_{1} A_{2}\right]$ et $\left[A_{2} A_{3}\right]$.
	
\begin{center}
\psset{xunit=1cm,yunit=1cm,arrowsize=4pt,plotpoints=1000,algebraic=true}
\begin{pspicture}(0,-1.1)(6.5,1.2)
%\psaxes[Dx = 1,Dy =1,labelsep=0.8mm,linewidth=.75pt,ticksize=-2pt 2pt,showorigin=false]{->}(0,0)(-0.5,-0.5)(8,8)
\psline(0,0)(6.5,0)
\pscircle(1,0){1}
\pscircle(3,0){1}
\pscircle(5,0){1}
\def\F{0.204*x}
\psplot{0}{6.5}{\F}
\uput[d](1,0){$O_{1}$}
\uput[d](3,0){$O_{2}$}
\uput[d](5,0){$O_{3}$}
\uput[d](1,-1){$\Gamma_{1}$}
\uput[d](3,-1){$\Gamma_{2}$}
\uput[d](5,-1){$\Gamma_{3}$}
\uput[r](6.5,0){$\Delta$}
\psdot[dotstyle=|](1,0)
\psdot[dotstyle=|](3,0)
\psdot[dotstyle=|](5,0)
\uput[l](0,0){$A$}
%\uput[dr](2.1,0.43){$B$}
%\uput[u](3.66,0.75){$C$}
%\uput[u](4.8,0.98){$D$}
\rput(2,0.6){$B$}
\rput(3.7,0.95){$C$}
\rput(4.75,1.18){$D$}
\end{pspicture}
\end{center}

	\n Par le point $A$, on mène une droite $(A D)$ tangente en $D$ au cercle $\Gamma_{3}$.
	\begin{enumerate}
		\item Montrer que la droite $(A D)$ coupe le cercle $\Gamma_{2}$ en deux points distincts $B$ et $C$. Calculer la longueur $B C$.
		\item Montrer que les droites $\left(B A_{1}\right)$ et $\left(C A_{2}\right)$ sont sécantes; on note $P$ leur point d'intersection.\\
		Montrer de même que les droites $\left(C A_{1}\right)$ et $\left(B A_{2}\right)$ sont sécantes; on note $Q$ leur point d'intersection.\\
		Que peut-on dire de la direction de la droite $(P Q)$ ?
	\end{enumerate}
	\item \textbf{Plus généralement}, on considère un entier $n$ strictement supérieur à 1 et $n$ cercles $\Gamma_{1}, \Gamma_{2}, \ldots, \Gamma_{n}$ de même rayon $r$ strictement positif, de centres respectifs $O_{1}, O_{2}, \ldots, O_{n}$ distincts et alignés dans cet ordre sur la demi-droite $\Delta$. Le cercle $\Gamma_{1}$ passe par $A$ et, pour tout $k>1$, le cercle $\Gamma_{k}$ est tangent au cercle $\Gamma_{k-1}$.\\
	Les diamètres des cercles $\Gamma_{1}, \Gamma_{2}, \ldots, \Gamma_{n}$ sur la droite $\Delta$ sont notés respectivement $\left[A A_{1}\right],\left[A_{1} A_{2}\right], \ldots,\left[A_{n-1} A_{n}\right]$.

\begin{center}
\psset{xunit=1cm,yunit=1cm,arrowsize=4pt,plotpoints=1000,algebraic=true}
\begin{pspicture}(0,-1.1)(10.5,1.2)
%\psaxes[Dx = 1,Dy =1,labelsep=0.8mm,linewidth=.75pt,ticksize=-2pt 2pt,showorigin=false]{->}(0,0)(-0.5,-0.5)(8,8)
\psline(0,0)(4,0)
\psline[linestyle=dashed](4,0)(5,0)
\psline(5,0)(7,0)
\psline[linestyle=dashed](7,0)(8,0)
\psline(8,0)(10.5,0)
\pscircle(1,0){1}
\pscircle(3,0){1}
\pscircle(6,0){1}
\pscircle(9,0){1}
\def\F{0.112*x}
\psplot{0}{3.9}{\F}
\psplot[linestyle=dashed]{3.9}{5.185}{\F}
\psplot{5.185}{6.667}{\F}
\psplot[linestyle=dashed]{6.667}{8.889}{\F}
\psplot{8.889}{10.5}{\F}
\uput[d](1,0){$O_{1}$}
\uput[d](3,0){$O_{2}$}
\uput[d](6,0){$O_{k}$}
\uput[d](9,0){$O_{n}$}
\uput[d](1,-1){$\Gamma_{1}$}
\uput[d](3,-1){$\Gamma_{2}$}
\uput[d](6,-1){$\Gamma_{k}$}
\uput[d](9,-1){$\Gamma_{n}$}
\uput[r](10.5,0){$\Delta$}
\psdot[dotstyle=|](1,0)
\psdot[dotstyle=|](3,0)
\psdot[dotstyle=|](6,0)
\psdot[dotstyle=|](9,0)
\uput[l](0,0){$A$}
%\uput[dr](1.975,0.221){$C_1$}
%\uput[u](2.026,0.227){$B_2$}
%\uput[u](3.9,0.436){$C_2$}
%\uput[u](5.185,0.58){$B_k$}
%\uput[u](6.667,0.745){$C_k$}
%\uput[u](8.889,0.994){$D$}


\rput(1.675,0.4){$C_1$}
\rput(2.35,0.45){$B_2$}
\rput(4.05,0.725){$C_2$}
\rput(5.05,0.85){$B_k$}
\rput(6.75,1.02){$C_k$}
\rput(8.889,1.225){$D$}
%\rput(2,0.6){$B$}
%\rput(3.7,0.95){$C$}
%\rput(4.75,1.18){$D$}
\end{pspicture}
\end{center}
	
	\n Par le point $A$, on mène une droite $(A D)$ tangente en $D$ au cercle $\Gamma_{n}$. Montrer que, pour tout $k$ entier tel que $1 \leq k \leq n-1$, cette droite coupe le cercle $\Gamma_{k}$ en deux points distincts $B_{k}$ et $C_{k}$ (on remarque que $B_{1}=A$ ).
	\newpage
	\begin{enumerate}
		\item Calculer la longueur $B_{k} C_{k}$ en fonction de $n$, de $k$ et de $r$.
		\begin{center}
		\textbf{Dans toute la suite du problème, on prend $r=1 .$ On pose $L(n, k)=B_{k} C_{k}$.}
		\end{center}
		\item Montrer que pour que $L(n, k)$ soit rationnel il faut et il suffit que la condition suivante soit vérifiée :
		\begin{center}
		$\left(\mathcal{C}_{1}\right)$\quad il existe $a \in \mathbb{N}$ tel que $n(n-1)-k(k-1)=4 a^{2}$
		\end{center}
	\end{enumerate}
\end{enumerate}


\section{Étude d'une surface}

\n L'espace est rapporté à un repère orthonormé $(O ; \vec{i}, \vec{j}, \vec{k})$. Les coordonnées (respectivement l'abscisse, l'ordonnée et la cote) d'un point sont notées $x, y$ et $z$.
On considère l'ensemble $\Sigma$ des points $M$ de coordonnées $(x, y, z)$ vérifiant :
$$z^{2}=x(x-1)-y(y-1)$$

\begin{enumerate}
	\item Soient $\lambda$ un réel et $P_{\lambda}$ le plan d'équation $x=\lambda$.\\
	Montrer que l'intersection de $\Sigma$ et de $P_{\lambda}$ est un cercle $C_{\lambda}$ dont on déterminera, en fonction de $\lambda$, le centre et le rayon.
	\item Soit $I$ le point de coordonnées $\left(\frac{1}{2}, \frac{1}{2}, 0\right)$ et $(d)$ la droite passant par $I$ de vecteur directeur $\vec{i}$.\\
	Montrer que la droite $(d)$ est un axe de symétrie de $\Sigma$. Déterminer et dessiner l'intersection de $\Sigma$ et du plan d'équation $y=\frac{1}{2}$.
	\item Reconnaître la nature de l'ensemble $\Sigma$.
	\item Soit un entier $n$ strictement supérieur à 2 et un entier $k$ tel que $1 \leq k \leq n-1$. Montrer que $L(n, k)$ est rationnel si, et seulement si, les points de $\Sigma$ d'abscisse $n$ et d'ordonnée $k$ ont pour cote un nombre entier pair.
\end{enumerate}

\newpage

\section{Étude d'une limite}

\n \textit{À partir de la configuration étudiée partie 1, \textbf{2.}, on définit $\lambda_{n}$ comme la proportion du segment $[A D]$ située à l'intérieur des cercles $\left(\Gamma_{k}\right)$, pour $1 \leq k \leq n-1$. Ainsi, on a $\lambda_{n}=\frac{1}{A D} \Sum{k=1}{n-1} B_{k} C_{k}$.}

\begin{enumerate}
	\item \textbf{Calculs d'intégrales}\\
	On définit la fonction $f$, de $[0,1]$ dans $\mathbb{R}$, par : pour tout $x \in[0,1], f(x)=\Int{0}{x}{\sqrt{1-t^{2}}}{t}$ puis la fonction $F$, de $\left[0, \frac{\pi}{2}\right]$ dans $\mathbb{R}$, par : pour tout $x \in\left[0, \frac{\pi}{2}\right]$, $F(x)=f(\sin x)$.
	\begin{enumerate}
		\item Montrer que la fonction $F$ est dérivable et calculer sa dérivée, notée $F^{\prime}$.
		\item Montrer que, pour tout $x$ dans $\left[0, \frac{\pi}{2}\right]$, $F(x)=\Int{0}{x}{\cos^2 t}{t}$.
		\item Sans chercher à calculer les intégrales, démontrer l'égalité $\Int{0}{\frac{\pi}{2}}{\cos ^{2}(t)}{t}=\Int{0}{\frac{\pi}{2}}{\sin ^{2}(t)}{t}$ et en déduire la valeur commune des deux intégrales.
		\item En déduire que $\Int{0}{1}{\sqrt{1-t^{2}}}{t}=\frac{\pi}{4}$; interpréter géométriquement ce résultat.
	\end{enumerate}
	\item \begin{enumerate}
		\item Montrer que pour tout $n \geq 2$, $\lambda_{n}=\frac{2}{2 n-1} \Sum{k=1}{n} \sqrt{1-\frac{k}{n} \cdot \frac{k-1}{n-1}}$.
		\item Montrer que si $n \geq 2$ et $1 \leq k \leq n$, on a : $\left(\frac{k-1}{n}\right)^{2} \leq \frac{k}{n} \cdot \frac{k-1}{n-1} \leq\left(\frac{k}{n}\right)^{2}$.
		\item On pose $I_{n, k}=\Int{\frac{k-1}{n}}{\frac{k}{n}}{\sqrt{1-t^{2}}}{t}$.\\
		Montrer que pour des valeurs convenables de $n$ et $k$, que l'on précisera, on a :
		$$n I_{n, k+1} \leq \sqrt{1-\frac{k-1}{n-1} \cdot \frac{k}{n}} \leq n I_{n, k-1}$$
	\end{enumerate}
	\item Démontrer, à partir des résultats des questions \textbf{1.} et \textbf{2.} ci-dessus, que la suite $\left(\lambda_{n}\right)$ est convergente et calculer sa limite.
\end{enumerate}

\newpage

\section{Étude de la condition $\left(\mathcal{C}_{1}\right)$}

\n On considère deux entiers $n$ et $k$ tels que $1 \leq k \leq n-1$.

\begin{enumerate}
	\item On pose $p=2 n-1$ et $q=2 k-1$. Montrer que le couple $(n, k)$ vérifie la condition $\mathcal{C}_{1}$ si, et seulement si, $(p, q)$ est un couple d'entiers naturels impairs tels que $q<p$, vérifiant la condition $\left(\mathcal{C}_{2}\right)$ suivante :
	\begin{center}
	$\left(\mathcal{C}_{2}\right)$\quad il existe $a \in \mathbb{N}$ tel que $p^{2}-q^{2}=16 a^{2}$
	\end{center}
	\item Soit $(p, q)$ un couple de nombres entiers naturels, tel qu'il existe deux entiers $u>0$ et $v>0$, de parités différentes, pour lesquels $p=u^{2}+v^{2}$ et $q=u^{2}-v^{2}$. Montrer que $(p, q)$ est un couple d'entiers naturels impairs tels que $q<p$ vérifiant la condition $\left(\mathcal{C}_{2}\right)$.
	\item On considère un couple $(p, q)$, d'entiers naturels impairs et premiers entre eux, tels que $q<p$ et vérifiant la condition $\left(\mathcal{C}_{2}\right)$. Montrer qu'il existe deux entiers naturels $u$ et $v$ de parités différentes tels que $p=u^{2}+v^{2}$ et $q=u^{2}-v^{2} .$ Calculer alors, en fonction de $u$ et de $v$, la valeur de l'entier $a$ qui intervient dans la condition $\left(\mathcal{C}_{2}\right)$.
\end{enumerate}

\newpage

\section{Nombre premier somme de deux carrés}

\n \textit{On se propose dans cette partie de déterminer tous les nombres premiers qui peuvent s'écrire comme somme de deux carrés d'entiers naturels. On désignera plus simplement un tel nombre comme étant << somme de deux carrés >>.}

\begin{enumerate}
	\item \begin{enumerate}
		\item Montrer que si $n$ est un entier naturel impair somme de deux carrés, il est congru à 1 modulo 4.
		\item Écrire 2 et 5 comme somme de deux carrés.
	\end{enumerate}
\end{enumerate}

\n \textit{Dans la suite de la partie 5, $p$ désigne un nombre premier congru à 1 modulo 4 et strictement supérieur à 5. On l'écrit sous la forme $p=4 m+1$ (avec $m>2$).\\
On définit $S=\left\{(x, y, z) \in \mathbb{N} \times \mathbb{N} \times \mathbb{Z} \mid 4 x y+z^{2}=p\right\}$.}

\begin{enumerate}\setcounter{enumi}{1}
	\item \begin{enumerate}
		\item Montrer que $S$ est un ensemble fini non vide et que l'intersection de $S$ et de l'ensemble d'équation $x=y+z$ est vide.
		\item À tout triplet $(x, y, z)$ de $S$, on associe le triplet $\left(x^{\prime}, y^{\prime}, z^{\prime}\right)$ défini par :
		$$\left(x^{\prime}, y^{\prime}, z^{\prime}\right)=\left\{\begin{array}{ll}
(x-y-z, y, 2 y+z) & \text { si } x>y+z \\
(y+z-x, x, 2 x-z) & \text { si } x<y+z
\end{array}\right.$$
		Montrer que pour tout $(x, y, z)$ de $S$, $\left(x^{\prime}, y^{\prime}, z^{\prime}\right)$ est aussi élément de $S$.
	\end{enumerate}
\end{enumerate}

\n \textit{On considère désormais la suite de triplets dans $S$ définie en itérant le procédé précédent de la manière suivante :
\begin{itemize}
	\item On part du triplet $\left(x_{0}, y_{0}, z_{0}\right)=(m, 1,1)$;
	\item $\left(x_{k}, y_{k}, z_{k}\right)$ ayant été défini dans $S$, on prend $x_{k+1}=x_{k}^{\prime}$, $y_{k+1}=y_{k}^{\prime}$, $z_{k+1}=z_{k}^{\prime}$.
\end{itemize}}

\begin{enumerate}\setcounter{enumi}{2}
	\item \begin{enumerate}
		\item \textit{Étude d'un cas particulier.} Dans cette question seulement, on prend $m=10$. Déterminer les triplets $\left(x_{k}, y_{k}, z_{k}\right)$ pour $0 \leq k \leq 11$.
		\item Montrer que si $(a, b, c)=\left(x_{k}, y_{k}, z_{k}\right)$, avec $k \geq 2$, alors le triplet $\left(x_{k-1}, y_{k-1}, z_{k-1}\right)$ est :
		$$\left\{\begin{array}{ll}
(a-b+c, b, c-2 b) & \text { si } a-4 b+2 c>0 \\
(b, a-b+c, 2 b-c) & \text { si } a-4 b+2 c<0
\end{array}\right.$$
		Montrer que ce résultat est encore vrai pour $k=1$.
		\item Montrer qu'il existe deux entiers distincts $k$ et $\ell$ tels que $\left(x_{k}, y_{k}, z_{k}\right)=\left(x_{\ell}, y_{\ell}, z_{\ell}\right)$.\\
		En déduire qu'il existe un entier $n$ strictement positif tel que $\left(x_{n}, y_{n}, z_{n}\right)=(m, 1,1)$.
	\end{enumerate}
\end{enumerate}

\n \textit{On note désormais $n$ le plus petit entier strictement positif tel que :}
$$\left(x_{n}, y_{n}, z_{n}\right)=(m, 1,1)$$

\begin{enumerate}\setcounter{enumi}{3}	
	\item \begin{enumerate}
		\item Calculer $\left(x_{n-1}, y_{n-1}, z_{n-1}\right)$ et $\left(x_{n-2}, y_{n-2}, z_{n-2}\right)$.
		\item Montrer que, pour tout entier $j$ tel que $1 \leq j \leq n$ :
		$$\left(x_{j-1}, y_{j-1}, z_{j-1}\right)=\left\{\begin{array}{ll}
\left(x_{n-j}, y_{n-j},-z_{n-j}\right) & \text { si } x_{j-1}>y_{j-1}+z_{j-1} \\
\left(y_{n-j}, x_{n-j}, z_{n-j}\right) & \text { si } x_{j-1}<y_{j-1}+z_{j-1}
\end{array}\right.$$
		\item Montrer que $n$ est impair. On pose désormais $n=2 r+1$.
		\item Montrer que $x_{r}=y_{r}$. En déduire qu'il existe une décomposition de $p$ en somme de deux carrés.
	\end{enumerate}
	\item \begin{enumerate}
		\item Déduire des questions précédentes un algorithme permettant de décomposer $p$ en somme de deux carrés.
		\item Donner le plus petit nombre premier supérieur à 40 qui est somme de deux carrés et, à l'aide de cet algorithme, en préciser une décomposition (on indiquera les triplets calculés aux différentes étapes de l'itération).
	\end{enumerate}
\end{enumerate}


\section{Retour au problème initial}


\begin{enumerate}
	\item Soient $n$ et $m$ deux entiers naturels somme de deux carrés, $n=a^{2}+b^{2}$, $m=c^{2}+d^{2}$. En introduisant les nombres complexes $a+i b$ et $c+i d$ et en considérant $n=|a+i b|^{2}$ et $m=|c+i d|^{2}$, montrer que le produit $m n$ est un entier somme de deux carrés et en donner explicitement une décomposition en fonction de $a$, $b$, $c$ et $d$.
	\item On se propose de démontrer, pour tout entier $n$ strictement positif, la proposition $\mathcal{P}(n)$ suivante : 
	\begin{center}
	$\mathcal{P}(n)$ : << tout nombre premier qui divise $n^{2}+1$ est somme de deux carrés >>.
	\end{center}
	Pour cela, on procède par récurrence sur $n$.
	\begin{enumerate}
		\item Montrer que $\mathcal{P}(1)$, $\mathcal{P}(2)$ et $\mathcal{P}(3)$ sont vraies.
		\item Soit $n$ un entier strictement supérieur à 1. On suppose la proposition $\mathcal{P}(j)$ vraie pour tout entier $j$ tel que $1 \leq j \leq n-1$ et on considère un nombre premier $p$ qui divise $n^{2}+1$.
		\begin{enumerate}
			\item Montrer que $p$ est différent de $n$.
			\item On suppose $p<n$. Montrer que $p$ divise $(n-p)^{2}+1$.
			\item On suppose $p>n$ et $p<n^{2}+1$. Montrer que les autres diviseurs premiers de $n^{2}+1$ sont strictement inférieurs à $n$. En déduire, en discutant selon la parité de $n$, que $p$ est congru à 1 modulo 4.
			\item Montrer que $p$ est somme de deux carrés.
		\end{enumerate}
		\item Conclure.
	\end{enumerate}
	\item \begin{enumerate}
		\item Pour $s$ entier supérieur ou égal à 2, on note $p_{s}$ le plus petit diviseur premier du nombre $(s !)^{2}+1$.\\
		Montrer que $p_{s}>s$ et que $p_{s}$ est somme de deux carrés.
		\item En déduire qu'il existe une infinité de nombres premiers somme de deux carrés.
	\end{enumerate}
	\item \begin{enumerate}
		\item Montrer qu'il existe une infinité de couples d'entiers $(n, k)$ avec $1 \leq k<n$ tels que $L(n, k)$ soit rationnel.
		\item Déterminer un entier $n$ tel qu'il existe plusieurs valeurs de $k$ pour lesquelles $L(n, k)$ est rationnel.
	\end{enumerate}
\end{enumerate}


\end{document}